%% kapa: exam=1 round=20 year=2009 period=1 form=A subject=economics subject_ko=경제원론 seq=22 number=62
%% source: Q-Net 감정평가사 기출문제 (cst00309) artlSeq=5024178 file=extracted/1008/5024178/2009년 제20회 감정평가사 1차시험 기출문제/1교시.hwp
%% conversion: kordoc 4.24.51 -> tools/gen1.py (auto); review=proofread-sonnet; flags=none
\begin{question}{62}
다음 ( )안에 알맞은 것은?
\begin{framedtext}
( )은 명목이자율에서 인플레이션율을 뺀 것으로, 예금의 경우 일정기간 동안 실질구매력의 상승률을 나타낸다.
\end{framedtext}
\begin{choices}
  \item 실질이자율
  \item 실질환율
  \item 물가상승률
  \item 경제성장률
  \item 실업률
\end{choices}
\end{question}
\answer{1}
